\section{La pila tecnológica de AI for Math: de la intuición a la demostración verificable}

Muchas personas reducen AI for Math a «si el modelo sabe resolver problemas de matemáticas». Esa es solo la capa más superficial. El objetivo real es que la IA parta de problemas, conjeturas y razonamiento informal, y avance paso a paso hacia enunciados formales, búsqueda de demostraciones, verificación por máquina y retroalimentación humana.

\begin{figure}[H]
\centering
\includegraphics[width=0.94\textwidth]{ai-for-math-stack.png}
\caption{La pila tecnológica de AI for Math: del problema a la verificación y, después, a la retroalimentación de los matemáticos.}
\end{figure}

\begin{knowledgebox}{Lectura de la figura: AI for Math no es una capacidad aislada del modelo}
Cada capa de la figura puede convertirse en un cuello de botella: comprender el problema requiere semántica y conocimiento de contexto; el razonamiento informal requiere intuición; la formalización requiere los lenguajes Lean/Coq; la búsqueda de demostraciones requiere una biblioteca de teoremas y tactics; la verificación requiere un proof assistant; por último, se necesita que los matemáticos juzguen la dirección y el significado del trabajo.
\end{knowledgebox}

\subsection{Términos clave: vocabulario de AI for Math}

\noindent\begin{tabularx}{\textwidth}{p{0.22\textwidth}p{0.32\textwidth}X}
\toprule
Término & Problema que resuelve & Relación con este episodio \\
\midrule
Lean & Lenguaje de matemáticas formales y proof assistant & Convierte demostraciones en lenguaje natural en demostraciones que una máquina puede comprobar. \\
Coq & Otro tipo de proof assistant & Muestra que la demostración formal no tiene a Lean como única vía. \\
Tactic & Paso de demostración automático o semiautomático & Permite que el modelo construya paso a paso demostraciones comprobables. \\
Theorem Proving & Demostración de teoremas & Una de las tareas centrales de AI for Math. \\
Verifier & Verificador & Determina si una demostración es válida y evita la alucinación de algo que «parece correcto». \\
Formalization & Formalización & Lleva artículos, definiciones y conjeturas a un sistema verificable. \\
\bottomrule
\end{tabularx}

\subsection{Por qué el proof assistant es decisivo}

Las demostraciones matemáticas en lenguaje natural suelen omitir pasos y dependen del contexto y la intuición del lector. Un proof assistant exige que cada paso cumpla las reglas. Para la IA, esto constituye una retroalimentación fuerte: la demostración generada o bien la acepta el kernel, o bien falla y devuelve el estado del objetivo. Esta retroalimentación es más rigurosa que una calificación subjetiva y resulta más adecuada para el entrenamiento y la evaluación.

\begin{importantbox}{La verificabilidad es la ventaja defensiva de AI for Math}
Las matemáticas son uno de los pocos campos de alto valor capaces de ofrecer señales de verificación fuertes. Sistemas como Lean/Coq hacen que la salida del modelo no solo «se parezca a una demostración», sino que deba superar una comprobación por máquina. Esto le da a AI for Math, al mismo tiempo, valor científico y capacidad de entrenarse desde la ingeniería.
\end{importantbox}

\subsection{Resumen de la sección}

La pila tecnológica de AI for Math abarca la intuición, los lenguajes formales, la búsqueda de demostraciones y la verificación por máquina. La verdadera dificultad está en conectar estas capas, no en resolver problemas en una sola de ellas.
